\documentclass[conference]{IEEEtran}
\IEEEoverridecommandlockouts
\usepackage{cite}
\usepackage{amsmath,amssymb,amsfonts}
\usepackage{algorithmic}
\usepackage{graphicx}
\usepackage{textcomp}
\usepackage{xcolor}
\usepackage{booktabs}
\usepackage{multirow}
\usepackage{url}
\def\BibTeX{{\rm B\kern-.05em{\sc i\kern-.025em b}\kern-.08em
    T\kern-.1667em\lower.7ex\hbox{E}\kern-.125emX}}

\begin{document}

\title{Does Temporal Scale Matter? A Comparative Study of Single-Scale vs.\ Multi-Scale Audio Feature Fusion for Music Information Retrieval}

\author{\IEEEauthorblockN{Gabriel Espinoza Hern\'andez}
\IEEEauthorblockA{\textit{Faculty of Computing} \\
\textit{Universidad de Ingenier\'ia y Tecnolog\'ia (UTEC)}\\
Lima, Per\'u \\
gabriel.espinoza.h@utec.edu.pe}
\and
\IEEEauthorblockN{Aurea Soriano-Vargas}
\IEEEauthorblockA{\textit{Faculty of Computing} \\
\textit{Universidad de Ingenier\'ia y Tecnolog\'ia (UTEC)}\\
Lima, Per\'u \\
0000-0002-8429-4119}
}

\maketitle

% ============================================================
\begin{abstract}
Handcrafted audio features such as MFCCs, chroma, and spectral descriptors remain widely used in Music Information Retrieval (MIR) pipelines.
However, the temporal window over which these features are computed is typically chosen by convention rather than by principled analysis.
This paper investigates how the temporal scale of feature extraction affects downstream MIR performance and whether combining features from multiple scales yields complementary information.
We extract a comprehensive set of audio descriptors at three temporal resolutions---short ($\sim$200\,ms), medium ($\sim$2\,s), and long ($\sim$5\,s)---and compare single-scale representations against an early-fusion multi-scale representation on genre classification (GTZAN, FMA-small) and auto-tagging (MagnaTagATune).
We further analyze feature importance at each scale to determine which descriptors benefit from longer or shorter contexts.
Our results are expected to show that multi-scale fusion consistently outperforms single-scale baselines, that timbral features (MFCCs, spectral centroid) are most discriminative at short scales while tonal and structural features (chroma, spectral contrast) benefit from longer windows, and that a simple concatenation strategy suffices to capture cross-scale complementarity without deep learning overhead.
\end{abstract}

\begin{IEEEkeywords}
Music Information Retrieval, Audio Features, Multi-Scale Analysis, Feature Fusion, Temporal Resolution
\end{IEEEkeywords}

% ============================================================
\section{Introduction}

Music audio analysis typically begins with the extraction of low-level descriptors---Mel-frequency cepstral coefficients (MFCCs), chroma vectors, spectral moments, and zero-crossing rates---computed over short overlapping frames of 20--50\,ms and then aggregated into summary statistics over a longer analysis window \cite{muller2015fundamentals, lerch2012introduction}.
The choice of this analysis window is crucial: too short and the feature misses the temporal context needed to characterize harmony or rhythm; too long and transient timbral details are averaged away.
Despite its importance, the effect of window size is seldom studied systematically; most pipelines simply adopt a conventional value (e.g., 1\,s or the length of the full clip) without justification.

Recent work in self-supervised music representation learning has highlighted the value of multi-scale temporal modeling.
Buisson et al.\ showed that learning embeddings at multiple structural levels improves segmentation \cite{buisson2024self}, while Alonso-Jim\'enez et al.\ demonstrated that tokenizing audio at a fixed temporal resolution produces strong but scale-limited representations \cite{alonso2025omar}.
These findings suggest that \emph{no single temporal scale captures all the musically relevant information}, but the evidence comes from complex deep models whose internal behavior is difficult to disentangle from the effect of scale itself.

This paper asks a simpler but foundational question: \textbf{does multi-scale feature extraction improve MIR performance even with classical handcrafted descriptors and standard machine-learning classifiers?}
Answering this question has both practical and theoretical value.
Practically, many MIR applications---especially in resource-constrained settings---still rely on handcrafted features \cite{ding2023audio, tamm2024comparative}.
Theoretically, isolating the effect of temporal scale from model architecture provides cleaner evidence about the nature of musical information at different time horizons.

Our contributions are:
\begin{enumerate}
    \item A systematic evaluation of how three temporal scales (short, medium, long) affect the discriminative power of seven widely used audio descriptors.
    \item An empirical comparison of single-scale representations against early-fusion multi-scale representations on genre classification and auto-tagging, using both traditional (Random Forest, SVM) and lightweight neural (MLP) classifiers.
    \item A feature-importance analysis that reveals which descriptors are most informative at each scale, providing actionable guidelines for MIR practitioners.
\end{enumerate}

% ============================================================
\section{Related Work}

\subsection{Audio Feature Extraction for MIR}

The standard MIR feature pipeline computes frame-level descriptors and aggregates them into track-level or segment-level statistics.
M\"uller's comprehensive treatment \cite{muller2015fundamentals} and Lerch's textbook \cite{lerch2012introduction} establish the canonical set: MFCCs for timbre, chroma for harmony, spectral centroid and rolloff for brightness, zero-crossing rate (ZCR) for noisiness, and spectral contrast for texture.
Despite decades of use, the temporal resolution at which these features are most effective has received surprisingly little dedicated study.
Peeters \cite{peeters2004large} proposed a large set of audio features with texture windows of 1--3\,s but did not compare across scales.
Lartillot et al.\ \cite{lartillot2008multi} introduced multi-scale analysis in MIRtoolbox but focused on rhythm and did not evaluate feature fusion for classification.

\subsection{Multi-Scale Representations in Music}

The importance of modeling music at multiple temporal levels has been recognized in structure analysis \cite{serra2014unsupervised} and, more recently, in deep learning.
Buisson et al.\ \cite{buisson2024self} learn continuous multi-level embeddings via contrastive objectives and demonstrate improvements in structural segmentation.
OMAR-RQ \cite{alonso2025omar} uses masked token prediction with vector quantization at a single temporal scale, achieving strong results across diverse tasks but without explicit multi-scale design.
SemiSupCon \cite{guinot2024semi} improves embedding geometry through semi-supervised contrastive learning but collapses temporal information into a global representation.
These deep approaches implicitly capture some multi-scale information through hierarchical layers, but the contribution of explicit temporal scale variation remains confounded with model capacity and training strategy.

\subsection{Feature Fusion Strategies}

Feature-level (early) fusion concatenates descriptors before classification, while decision-level (late) fusion combines classifier outputs \cite{hamel2011temporal}.
Hamel and Eck \cite{hamel2011temporal} showed that temporal pooling strategies significantly affect classification accuracy.
Li et al.\ \cite{li2003comparative} compared frame-level and clip-level features for genre classification.
More recently, Tamm and Aljanaki \cite{tamm2024comparative} demonstrated that different pretrained embeddings excel on different downstream tasks, indirectly suggesting that no single feature view is sufficient.
Ding and Lerch \cite{ding2023audio} showed that even simple student models can match teacher embeddings when transfer is done carefully, supporting the idea that feature engineering---including temporal design---still matters.

Our work differs from all of the above by explicitly isolating temporal scale as the independent variable across a controlled set of handcrafted features and standard classifiers, providing evidence that is not confounded by deep model architecture or training procedure.

% ============================================================
\section{Research Questions and Hypotheses}

\textbf{RQ1.} How does the temporal window size (short $\sim$200\,ms, medium $\sim$2\,s, long $\sim$5\,s) affect the discriminative quality of individual audio features for genre classification and auto-tagging?

\textbf{RQ2.} Does concatenating features extracted at multiple temporal scales (early fusion) outperform any single-scale representation on the same downstream tasks?

\textbf{RQ3.} Which audio descriptors are most important at each temporal scale, and does this importance shift across tasks?

\medskip
\noindent\textbf{H1.} Timbral features (MFCCs, spectral centroid, ZCR) will be most discriminative at short scales ($\sim$200\,ms), while tonal and structural features (chroma, spectral contrast, tonnetz) will benefit from longer windows ($\geq$2\,s).

\noindent\textbf{H2.} Multi-scale feature fusion will consistently outperform the best single-scale baseline across all tasks and classifiers, because features at different scales capture complementary aspects of musical content.

\noindent\textbf{H3.} Feature importance rankings will differ across scales, confirming that each scale emphasizes different musical properties.

% ============================================================
\section{Methodology}

\subsection{Overview}

The experimental pipeline consists of four stages (Fig.~\ref{fig:pipeline}): (1)~audio preprocessing, (2)~multi-scale feature extraction, (3)~feature aggregation and fusion, and (4)~classification and analysis.
All code will be implemented in Python using \texttt{librosa} \cite{mcfee2015librosa} for feature extraction, \texttt{scikit-learn} for classical classifiers, and \texttt{PyTorch} for the MLP baseline.
Experiments will run on Google Colab Pro.

\begin{figure}[t]
\centering
\fbox{\parbox{0.95\columnwidth}{\centering
\small
\textbf{Audio Waveform (16\,kHz mono)} \\[4pt]
$\downarrow$ \\[2pt]
\textbf{Windowing:} Short (200\,ms) $|$ Medium (2\,s) $|$ Long (5\,s) \\[4pt]
$\downarrow$ \\[2pt]
\textbf{Feature Extraction per window:} \\
MFCCs, Chroma, Spectral Centroid, Spectral Contrast, \\
Spectral Rolloff, ZCR, Tonnetz \\[4pt]
$\downarrow$ \\[2pt]
\textbf{Aggregation:} mean $\pm$ std per feature per scale \\[4pt]
$\downarrow$ \\[2pt]
\textbf{Representations:} \\
Single-scale ($\mathbf{x}_s$, $\mathbf{x}_m$, $\mathbf{x}_l$) ~~$|$~~ Multi-scale ($[\mathbf{x}_s \| \mathbf{x}_m \| \mathbf{x}_l]$) \\[4pt]
$\downarrow$ \\[2pt]
\textbf{Classifiers:} Random Forest, SVM, MLP \\[4pt]
$\downarrow$ \\[2pt]
\textbf{Evaluation \& Feature Importance Analysis}
}}
\caption{Overview of the multi-scale feature extraction and evaluation pipeline.}
\label{fig:pipeline}
\end{figure}

\subsection{Datasets}

We use three publicly available datasets that cover the two core tasks (Table~\ref{tab:datasets}):

\begin{table}[t]
\centering
\caption{Datasets used in this study.}
\label{tab:datasets}
\begin{tabular}{@{}lccc@{}}
\toprule
\textbf{Dataset} & \textbf{Task} & \textbf{Tracks} & \textbf{Classes/Tags} \\
\midrule
GTZAN \cite{tzanetakis2002musical} & Genre classif. & 1\,000 & 10 genres \\
FMA-small \cite{defferrard2017fma} & Genre classif. & 8\,000 & 8 genres \\
MagnaTagATune \cite{law2009evaluation} & Auto-tagging & 25\,863 & 50 tags \\
\bottomrule
\end{tabular}
\end{table}

\begin{itemize}
    \item \textbf{GTZAN} \cite{tzanetakis2002musical}: 1\,000 tracks of 30\,s each, 10 genres. Despite known limitations (repetitions, mislabelings), it remains a standard benchmark for feature-level studies. We use the fault-filtered split of Kereliuk et al.\ \cite{kereliuk2015counterexample}.
    \item \textbf{FMA-small} \cite{defferrard2017fma}: 8\,000 tracks of 30\,s each, balanced across 8 top-level genres. Provides a larger and cleaner alternative to GTZAN.
    \item \textbf{MagnaTagATune (MTAT)} \cite{law2009evaluation}: 25\,863 clips of $\sim$29\,s each, annotated with 50 binary tags covering genre, instrumentation, and mood. Multi-label setting tests whether multi-scale features help in a more complex output space.
\end{itemize}

\subsection{Audio Preprocessing}

All audio is resampled to 16\,kHz mono.
Silence at the beginning and end of each track is trimmed using a top-dB threshold of 20\,dB.
No further augmentation is applied in the main experiments to isolate the effect of temporal scale.

\subsection{Multi-Scale Feature Extraction}

For each track, we divide the audio into non-overlapping analysis windows at three scales:

\begin{itemize}
    \item \textbf{Short ($w_s = 200$\,ms):} Captures instantaneous timbral and spectral properties. A 30\,s track yields $\sim$150 windows.
    \item \textbf{Medium ($w_m = 2$\,s):} Captures short musical phrases, local harmonic context, and rhythmic patterns. Yields $\sim$15 windows per track.
    \item \textbf{Long ($w_l = 5$\,s):} Captures structural tendencies, sustained harmonic fields, and longer textural patterns. Yields $\sim$6 windows per track.
\end{itemize}

Within each analysis window, we compute frame-level features using \texttt{librosa} with a hop length of 512 samples (32\,ms at 16\,kHz) and an FFT size of 2048:

\begin{enumerate}
    \item \textbf{MFCCs} (20 coefficients): Compact representation of spectral envelope, widely used as a timbral descriptor.
    \item \textbf{Chroma} (12 bins): Pitch-class energy distribution, capturing harmonic and tonal content.
    \item \textbf{Spectral Centroid} (1 value): Center of mass of the spectrum, correlating with perceived brightness.
    \item \textbf{Spectral Contrast} (7 bands): Difference between peaks and valleys in the spectrum across frequency bands.
    \item \textbf{Spectral Rolloff} (1 value): Frequency below which a specified percentage (85\%) of total spectral energy is contained.
    \item \textbf{Zero-Crossing Rate} (1 value): Rate at which the signal changes sign, related to noisiness and percussiveness.
    \item \textbf{Tonnetz} (6 values): Tonal centroid features representing harmonic relations in the Tonnetz space.
\end{enumerate}

This yields $d = 20 + 12 + 1 + 7 + 1 + 1 + 6 = 48$ frame-level features per frame.

\subsection{Feature Aggregation}

For each analysis window, frame-level features are aggregated into a fixed-length vector by computing the \textbf{mean} and \textbf{standard deviation} across frames, yielding $2 \times 48 = 96$ dimensions per window.
The track-level representation at scale $k \in \{s, m, l\}$ is then obtained by computing the mean and standard deviation of the window-level vectors across all windows of that scale, yielding $\mathbf{x}_k \in \mathbb{R}^{192}$.

The four representations compared are:
\begin{align}
    \text{Short-only:} \quad & \mathbf{x}_s \in \mathbb{R}^{192} \\
    \text{Medium-only:} \quad & \mathbf{x}_m \in \mathbb{R}^{192} \\
    \text{Long-only:} \quad & \mathbf{x}_l \in \mathbb{R}^{192} \\
    \text{Multi-scale:} \quad & \mathbf{x}_{\text{ms}} = [\mathbf{x}_s \| \mathbf{x}_m \| \mathbf{x}_l] \in \mathbb{R}^{576}
\end{align}

All features are standardized (zero mean, unit variance) using statistics computed only on the training set.

\subsection{Classification Models}

To ensure that observed differences are attributable to the features rather than model capacity, we use three classifiers of increasing complexity:

\begin{itemize}
    \item \textbf{Random Forest (RF):} 500 trees, max depth 30. Provides built-in feature importance via mean decrease in impurity (MDI).
    \item \textbf{Support Vector Machine (SVM):} RBF kernel, hyperparameters ($C$, $\gamma$) tuned via 3-fold cross-validation on the training set.
    \item \textbf{Multi-Layer Perceptron (MLP):} Two hidden layers (256, 128), ReLU activation, dropout 0.3, trained for 100 epochs with Adam ($\text{lr} = 10^{-3}$) and early stopping (patience 10).
\end{itemize}

For GTZAN: stratified 10-fold cross-validation.
For FMA-small: official train/validation/test splits.
For MTAT: standard 12:1:3 split.

\subsection{Evaluation Metrics}

\begin{itemize}
    \item \textbf{Genre classification} (GTZAN, FMA-small): Overall accuracy, macro-averaged F1 score, and per-class F1.
    \item \textbf{Auto-tagging} (MTAT): Mean average precision (mAP), macro-averaged ROC-AUC, and per-tag AP for the top-10 most frequent tags.
\end{itemize}

\subsection{Feature Importance Analysis (RQ3)}

To answer RQ3, we analyze feature importance using three complementary methods:

\begin{enumerate}
    \item \textbf{Mean Decrease in Impurity (MDI):} From Random Forest, averaged across folds/splits.
    \item \textbf{Permutation Importance:} Measure accuracy drop when each feature is randomly shuffled, applied to the best-performing classifier per task.
    \item \textbf{Scale-Conditional Importance:} For the multi-scale representation, we group features by their originating scale ($s$, $m$, $l$) and compute the aggregate importance of each scale. We further compute the importance of each descriptor type (MFCCs, chroma, etc.) within each scale, yielding a $7 \times 3$ importance matrix that reveals which descriptors benefit most from which temporal resolution.
\end{enumerate}

\subsection{Statistical Validation}

All comparisons between single-scale and multi-scale representations are tested for statistical significance using the Wilcoxon signed-rank test ($\alpha = 0.05$) applied to per-fold accuracy (GTZAN) or per-tag AP (MTAT).
Effect sizes are reported using Cliff's delta.

% ============================================================
\section{Experimental Plan and Timeline}

Given the constraint of one month and access to Google Colab Pro, we organize the work into four weekly sprints:

\begin{table}[t]
\centering
\caption{Proposed 4-week timeline.}
\label{tab:timeline}
\begin{tabular}{@{}cp{6.2cm}@{}}
\toprule
\textbf{Week} & \textbf{Tasks} \\
\midrule
1 & Dataset preparation (download GTZAN and MTAT; verify FMA-small). Implement feature extraction pipeline at three scales. Extract and cache all features to disk. \\
\midrule
2 & Train and evaluate single-scale classifiers (RF, SVM, MLP) on all three datasets. Record per-fold/per-split metrics. \\
\midrule
3 & Train and evaluate multi-scale fusion classifiers. Run feature importance analysis. Statistical tests. Ablation: pairwise scale combinations ($s{+}m$, $s{+}l$, $m{+}l$). \\
\midrule
4 & Write results and discussion. Generate figures and tables. Final paper compilation and review. \\
\bottomrule
\end{tabular}
\end{table}

\subsection{Computational Requirements}

Feature extraction with \texttt{librosa} is CPU-bound and does not require GPU.
For the three datasets combined ($\sim$35\,000 tracks $\times$ 30\,s), extraction at all three scales is estimated at $\sim$8--12 hours of CPU time on Colab Pro.
Features will be cached as NumPy arrays ($\sim$150\,MB total) to allow rapid iteration on classifiers.
RF and SVM training require only CPU.
MLP training on $\sim$35\,000 samples with 576 features is lightweight and completes in minutes per configuration.

% ============================================================
\section{Expected Results}

We anticipate the following outcomes:

\begin{enumerate}
    \item \textbf{Scale-dependent feature quality (RQ1):} MFCCs and spectral centroid will achieve their highest single-scale accuracy at the short scale, because timbral information is captured in instantaneous spectral shape. Chroma and tonnetz will peak at the medium scale, where harmonic progressions become visible. Spectral contrast will benefit from longer windows that expose sustained textural differences between genres.

    \item \textbf{Multi-scale advantage (RQ2):} The concatenated multi-scale representation will outperform every single-scale baseline by 2--5 percentage points in accuracy (genre classification) and 1--3 points in mAP (auto-tagging), across all classifiers. The improvement will be largest for tasks with diverse label semantics (MTAT), where tags span timbral, harmonic, and structural attributes.

    \item \textbf{Feature importance shifts (RQ3):} The $7 \times 3$ importance matrix will show a clear diagonal tendency---short-scale timbral features dominating for percussion/instrument tags, medium-scale tonal features for genre and mood, and long-scale contrast features for structural tags. This will validate the hypothesis that different musical properties reside at different temporal resolutions.
\end{enumerate}

% ============================================================
\section{Discussion Directions}

The expected results open several discussion threads:

\begin{itemize}
    \item \textbf{Implications for deep models:} If simple handcrafted features already benefit from multi-scale extraction, this provides empirical grounding for the architectural choices in recent deep multi-scale models \cite{buisson2024self, alonso2025omar}---their advantage may partly stem from implicit multi-scale processing rather than from learned representations per se.

    \item \textbf{Practical guidelines:} Many production MIR systems still use handcrafted features for efficiency \cite{tamm2024comparative}. Our analysis will provide concrete recommendations on optimal window sizes for each descriptor and task.

    \item \textbf{Limitations and future work:} The study uses non-overlapping windows and simple concatenation. Future work could explore overlapping windows, attention-based fusion, and the interaction between multi-scale handcrafted features and pretrained deep embeddings (e.g., using multi-scale features as complementary inputs to models like OMAR-RQ \cite{alonso2025omar} or SemiSupCon \cite{guinot2024semi}).
\end{itemize}

% ============================================================
\section{Conclusion}

This proposal describes a controlled empirical study to determine whether and how temporal scale affects the quality of handcrafted audio features for MIR.
By systematically comparing single-scale and multi-scale representations on genre classification and auto-tagging with standard classifiers, we aim to provide foundational evidence that multi-scale feature fusion captures complementary musical information across timbral, harmonic, and structural dimensions.
The study is designed to be reproducible, computationally lightweight, and directly useful to both researchers designing new representation learning architectures and practitioners building MIR systems under resource constraints.

\bibliographystyle{IEEEtran}
\bibliography{references}

\end{document}
